\documentclass[a4paper,11pt]{article}
\usepackage[T1]{fontenc}
\usepackage[utf8]{inputenc}
\usepackage{lmodern,microtype}
\usepackage[a4paper,margin=1in,headheight=15pt,headsep=14pt]{geometry}
\usepackage{amsmath,amssymb,graphicx,booktabs}
\usepackage[table,svgnames]{xcolor}
\usepackage[most]{tcolorbox}
\tcbuselibrary{skins,breakable}
\usepackage{pifont}
\IfFileExists{bbding.sty}{\usepackage{bbding}\newcommand{\glyphHand}{\Pointinghand}\newcommand{\glyphPencil}{\Pencil}}{\newcommand{\glyphHand}{\ding{43}}\newcommand{\glyphPencil}{\ding{46}}}
\usepackage{enumitem}
\setlist{leftmargin=1.5em,topsep=4pt,itemsep=3pt}
\newlist{steps}{enumerate}{1}
\setlist[steps]{label=\textbf{Step \arabic*.},leftmargin=4.4em,labelwidth=3.8em,labelsep=0.6em,labelindent=0pt,align=left,topsep=5pt,itemsep=6pt}
\usepackage{parskip,fancyhdr,titlesec}
\usepackage[hidelinks]{hyperref}
\usepackage{xurl}
\hypersetup{breaklinks=true}
\definecolor{navy}{HTML}{21355E}\definecolor{navyrule}{HTML}{21355E}\definecolor{inktext}{HTML}{1A202C}\definecolor{mutedtext}{HTML}{6B7280}
\definecolor{intuFrame}{HTML}{2C5AA0}\definecolor{intuFill}{HTML}{F4F7FC}
\definecolor{everyFrame}{HTML}{8A5A1E}\definecolor{everyFill}{HTML}{FBF1E3}
\definecolor{workFrame}{HTML}{2E7D52}\definecolor{workFill}{HTML}{F1F8F3}
\definecolor{watchFrame}{HTML}{B23A48}\definecolor{watchFill}{HTML}{FBF1F2}
\definecolor{keyFrame}{HTML}{6A4C93}\definecolor{keyFill}{HTML}{F4F0F9}
\color{inktext}
\hypersetup{linkcolor=navy,urlcolor=navy,citecolor=navy}
\newcommand{\CourseShort}{Machine Learning}
\newcommand{\SessionNo}{8}
\newcommand{\LectureTitle}{Advanced Decision Trees \& Pruning}
\pagestyle{fancy}\fancyhf{}
\fancyhead[L]{\footnotesize\color{mutedtext}\textsf{\CourseShort}}
\fancyhead[R]{\footnotesize\color{mutedtext}\textsf{Session \SessionNo\ \textperiodcentered\ \LectureTitle}}
\fancyfoot[L]{\footnotesize\color{mutedtext}\textsf{WILP, BITS Pilani}}
\fancyfoot[C]{\footnotesize\color{mutedtext}\thepage}
\renewcommand{\headrulewidth}{0.4pt}\renewcommand{\footrulewidth}{0pt}
\titleformat{\section}{\sffamily\Large\bfseries\color{navy}}{\thesection}{0.8em}{}[{\vspace{2pt}\color{navy!70}\titlerule[0.8pt]}]
\titleformat{\subsection}{\sffamily\large\bfseries\color{navy!92}}{\thesubsection}{0.7em}{}
\titlespacing*{\section}{0pt}{16pt}{8pt}\titlespacing*{\subsection}{0pt}{12pt}{4pt}
\newif\ifmlkfirstsec \mlkfirstsectrue
\newcommand{\sectionbreak}{\ifmlkfirstsec\mlkfirstsecfalse\else\clearpage\fi}
\newcommand{\secsub}[1]{\noindent{\sffamily\bfseries\itshape\color{navy!85}\large #1}\par\medskip}
\tcbset{housecallout/.style 2 args={enhanced, breakable, colback=#2, colframe=#1, boxrule=0.6pt, arc=2.4pt, outer arc=2.4pt, left=10pt, right=10pt, top=9pt, bottom=9pt, before skip=11pt, after skip=11pt, fontupper=\normalsize, coltitle=white, fonttitle=\bfseries\sffamily\small, attach boxed title to top left={xshift=6mm,yshift=-3mm}, boxed title style={colback=#1, colframe=#1, rounded corners, arc=3pt, boxrule=0pt, left=6pt, right=6pt, top=2pt, bottom=2pt}}}
\newtcolorbox{intuition}{housecallout={intuFrame}{intuFill}, title={\raisebox{-0.05ex}{\glyphHand}\,\ The intuition}}
\newtcolorbox{worked}[1]{housecallout={workFrame}{workFill}, title={\raisebox{-0.05ex}{\glyphPencil}\,\ Worked example: #1}}
\newtcolorbox{keytake}{housecallout={keyFrame}{keyFill}, title={\raisebox{-0.05ex}{\ding{51}}\,\ Key takeaway}}
\newcommand{\bitswordmark}{{\sffamily\bfseries\color{white}\fontsize{12.5}{15}\selectfont WILP, BITS Pilani}}
\newcommand{\titlebanner}[4]{\begin{tcolorbox}[enhanced, breakable=false, colback=navy, colframe=navy, boxrule=0pt, arc=3pt, outer arc=3pt, left=16pt, right=16pt, top=16pt, bottom=16pt, before skip=2pt, after skip=14pt, width=\linewidth]\noindent\begin{minipage}[c]{0.72\linewidth}\bitswordmark\end{minipage}\hfill\par\vspace{9pt}{\sffamily\footnotesize\bfseries\color{white!85}\MakeUppercase{#1}}\par\vspace{6pt}{\sffamily\bfseries\fontsize{26}{30}\selectfont\color{white} #2\par}\vspace{5pt}{\sffamily\large\color{white!88} #3\par}\vspace{9pt}{\color{white!55}\rule{\linewidth}{0.4pt}}\par\vspace{6pt}{\sffamily\footnotesize\color{white!82} #4\par}\end{tcolorbox}}
\newcommand{\housefig}[2]{\begin{figure}[htbp]\centering\includegraphics[width=\linewidth]{#1}\caption{#2}\end{figure}}
\newcommand{\vocab}[1]{\textbf{\color{navy}#1}}
\begin{document}
\thispagestyle{fancy}
\titlebanner{Machine Learning \textperiodcentered\ Companion Reader}{Advanced Decision Trees \& Pruning}{Continuous Attributes, Overfitting, Pre/Post-Pruning, Gini Index, Gain Ratio, and MDL}{Session 8 \textperiodcentered\ Read alongside the lecture slides \textperiodcentered\ Every slide example worked out in full}
\smallskip\noindent{\small\color{mutedtext}Prepared for the Work Integrated Learning Programmes (WILP), BITS Pilani.}\smallskip
\section{Continuous Attributes in Decision Trees}
\secsub{Continuous features are converted into binary decision thresholds by sorting values and testing candidate midpoints.}
Given sorted values $a_1 < a_2 < \dots < a_N$, candidate midpoints $c_i = (a_i + a_{i+1})/2$ are evaluated for maximum Information Gain.
\housefig{figures/fig_continuous_split.png}{Candidate split thresholds for a continuous attribute.}
\section{Model Overfitting \& Tree Pruning}
\secsub{Overfitting occurs when a tree memorizes training noise at the expense of generalization accuracy.}
\housefig{figures/fig_tree_overfitting.png}{Training versus test error as decision tree complexity increases.}
\subsection{Pruning Strategies}
\begin{enumerate}
  \item \textbf{Pre-Pruning:} Halt node expansion early if Information Gain is below a threshold.
  \item \textbf{Post-Pruning:} Trim subtrees from a fully grown tree using validation data.
\end{enumerate}
\section{Alternative Impurity Measures: Gini \& Gain Ratio}
\secsub{Gini Index and Gain Ratio address the limitations of standard Information Gain.}
\housefig{figures/fig_impurity_comparison.png}{Comparison of impurity metrics: Gini Index vs Scaled Entropy.}
\begin{equation}
\text{Gini}(S) = 1 - \sum_{i=1}^c p_i^2, \quad \text{GainRatio}(S, A) = \frac{\text{Gain}(S, A)}{\text{SplitInfo}(S, A)}.
\end{equation}
\section{Minimum Description Length (MDL) Principle}
\secsub{MDL balances model complexity against training error.}
\begin{equation}
\text{Cost}(\text{Model}, \text{Data}) = \text{Cost}(\text{Model}) + \text{Cost}(\text{Data} \mid \text{Model}).
\end{equation}
\section*{Glossary \& Symbols}
\noindent\textbf{Key Terms}
\begin{description}[leftmargin=8.5em,style=nextline,font=\normalfont\bfseries\color{navy}]
  \item[Pre-Pruning] Halting tree expansion early during construction.
  \item[Gini Index] An impurity metric measuring class overlap ($1 - \sum p_i^2$).
\end{description}
\end{document}
